\subsection{Blind Classification}
\label{sec:blind_classification}

\paragraph{Task Description.}
The blind classification experiment evaluates zero-shot topical categorisation
on the AG News dataset (\texttt{sh0416/ag\_news}) \emph{without} label
supervision.  Unlike the standard text classification variant, category
labels are never appended during ingestion---the system must discover
topical structure purely from value co-occurrence patterns in the article
text.  At test time the system must surface the correct category word
through associative recall alone.

\paragraph{Results.}
Table~\ref{tab:blind_classification_metrics} summarises the classification
metrics across $N = 100$ test samples.
The confusion matrix is shown in Figure~\ref{fig:blind_classification_confusion}.

\paragraph{Assessment.}
Blind classification accuracy was low.  Without label supervision the
substrate relies entirely on structural similarity between article content
and category words.  Scaling ingestion volume or enriching the corpus with
category-adjacent vocabulary may improve attractor formation.

\begin{table}[htbp]
  \centering
  \caption{Blind Classification --- summary metrics.}
  \label{tab:blind_classification_metrics}
  \begin{tabular}{ll}
    \toprule
    \textbf{Metric} & \textbf{Value} \\
    \midrule
    Overall Accuracy  & 11.0\% \\
    Balanced Accuracy & 21.7\% \\
    Macro-F1          & 0.227 \\
    Mean Resonance    & 0.1100 \\
    Sample Size       & $N = 100$ \\
    \bottomrule
  \end{tabular}
\end{table}
